\head{\emph{\textsc{Ilíada}}, Canto VI}
\section*{\centering\textbf{\textsc{Canto} VI}\footnote{\ehuno, pp. 42-43 + \ehdos, pp. 86, 94-97 + \ehctr, pp. 262-265.}}
\addcontentsline{toc}{subsection}{\textsc{Canto VI}}
%
\beginnumbering
\noindent-- -- -- -- -- --\par
\pstart
\setline{123}
¿Quién de entre los mortales eres tú el esforzado?\\
Jamás te ví en la lucha que a los héroes renombra;\\
mas, por tu audacia a todos los has sobrepasado,\\
con afrontar mi lanza, la de la larga sombra.\par
\pend
\noindent-- -- -- -- -- --\par
%
\pstart
\setline{128}
\firstlinenum{129}
Mas, si del cielo vienes, por ser un inmortal,\\
no pelearé, por cierto, contra un Dios celestial.\par
\pend
\noindent-- -- -- -- -- --\par
%
\pstart
\setline{142}
\firstlinenum{143}
Mas, si un mortal que nutre la tierra, eres por suerte,\\
acércate y más pronto te ultimará la muerte.\par
\pend
\noindent-- -- -- -- -- --\par
%
\pstart
\setline{145}
\firstlinenum{145}
Magnánimo Tideides, ¿a qué indagas mi origen,\\
si el del hombre es como el de las hojas pasajeras\\
que el viento en tierra esparce, cuando ya las otras se erigen\\
al pintar en el fértil bosque de la primavera?\\
Tal las generaciones nacen y acaban; pero\\
ya que saber desear nuestro bien conocido\\
linaje, que te enteres cumplidamente quiero.\par
\pend
\noindent-- -- -- -- -- --\par
%
\pstart
\setline{212}
Tal se expresó. Diomedes el ilustre en la guerra,\\
contento de escucharlo, clavó su lanza en tierra,\\
y así al pastor de pueblos dijo con voz amable:\par
--Eres mi antiguo huésped paterno, en verdad, pues\\
el deífico Eneo dió posada una vez\\
por veinte días en su palacio al intachable\\
Belerofonte, y ambos se obsequiaron cumplidos:\\
dió Eneo un cinto espléndido, de púrpura teñido,\\
y Belerofonte una copa de oro con asas,\\
que cuando yo me vine para acá, quedó en casa.\\
No recuerdo\footnote{En el impreso figura como «recuerbo», probablemente un error de imprenta.} a Tideo, puesto que me dejó\\
muy niño, cuando en Tebas pereció el pueblo aqueo;\\
mas, tu amistoso huesped ser en Argos deseo,\\
cual lo serás tú en Licia si a ese país voy yo.\\
Evitemos, pues, nuestras lanzas en la refriega,\\
que a mí muchos troyanos y aliados qué matar\\
me quedan todavía, si un dios me los allega,\\
y a tí muchos aqueos que podrás ultimar.\\
Y ahora cambiemos armas, y así éstos podrán ver\\
que huéspedes paternos nos gloriamos de ser.\par
\pend
\separator
%
\pstart
\setline{232}
Dicho esto, y de los carros descendiendo en seguida,\\
danse las manos para que su amistad se pruebe.\\
Entonces turbó a Glauco la razón Zeus Kronida;\\
pues al cambiar sus armas con Diomedes Tidida,\\
dió aquél oro por bronce, cien bueyes contra nueve.\par
\pend
\separator
%
\pstart
\setline{274}
Y prométele doce terneras, que inmoladas\\
de un año, y aun no uncidas, le serán, si se apiada\\
de la ciudad, esposas y chiquillos troyanos,\\
y al hijo de Tideo, guerrero que inhumano\\
nos amedrenta, aleja nuestra Ilión sagrada.\par
\pend
\noindent-- -- -- -- -- --\par
%
\pstart
\setline{305}
Divina entre las diosas, venerable Atenea,\\
Patrona nuestra, rompe la lanza de Diomedes\\
y haz que caiga delante de las Puertas Esceas.\par
\pend
\noindent-- -- -- -- -- --\par
%
\pstart
\setline{344}
Oh hermano de esta perra funesta y horrorosa,\\
¿por qué cuando mi madre me dió a luz, un mal viento\\
no me arrojó a los montes o al mar que, violento,\\
me habría arrastrado, antes de ocurrir estas cosas?\\
Pero, ya que los dioses tenían el intento\\
de causar tantos males, yo debí ser la esposa\\
de un hombre mejor, que ante la condena y la afrenta\\
de los otros varones, tomara su honra en cuenta.\\
De firmeza y corage siempre éste ha carecido,\\
hasta que un día, creo, tendrá su merecido.\\
Mas, entra hermano, y siéntate aquí, y reposa en calma,\\
del afán que por esta perra te angustia el alma,\\
y a causa de Alejandro; pues fatal anatema\\
nos ha infligido Zeus, para que en adelante,\\
sirva a los venideros nuestro baldón de tema.\par
\pend
\separator
%
\pstart
\setline{390}
Dijo el ama. Hérctor, yéndose de la casa muy presto,\\
por las bien hechas calles desanduvo el camino.\\
Ya en las puertas Esceas (que es de donde, al lado opuesto,\\
debía sobre el llano salir) corriendo vino\\
hacia él su esposa Andrómaca, opulenta heredera\\
del magnánimo Eécion que a la falda viviera\\
del Placo umbroso, en Tebas la hipoplaciana, donde\\
fué rey de los cilicios. A su hija poseía,\\
pues, Héctor del broncíneo casco. Así ella venía\\
ante él, y a la par suya marchaba una doncella,\\
llevando al seno el hijo de Héctor, un tierno infante,\\
que era su bienamado, lindo cual una estrella.\\
Escamandrio llamábalo él; mas los otros, como\\
sólo Héctor defendía la ciudad con aplomo,\\
Astianax le decían. El héroe, con cariño,\\
sonriendo en silencio contemplaba a su niño.\\
Andrómaca, llorosa, detúvose a su lado,\\
\dlgo{y de su mano asiéndole, díjole:}{--Infortunado,}\\
tu valor va a perderte, y en tanto no te apiadas\\
ni de tu hijo chiquito, ni de mí, desdichada,\\
que seré tu viuda pronto, pues sobre ti\\
caerán los aqueos para inmolarte así.\\
Privada de tu apoyo, morir será mi anhelo,\\
que tu muerte dejárame un dolor sin consuelo.\\
Ya mi padre no existe, mi augusta madre es muerta.\\
mató a mi padre Aquiles, asoló la poblada\\
ciudad de los cilicios, Tebas la de altas puertas,\\
mas no despojó a Fécion; antes bien, compungido,\\
quemólo con sus armas y un túmulo ha erigido,\\
que de álamos cercaron las ninfas Orcades\\
hijas de Zeus portaégida. Mis siete hermanos juntos,\\
en un día bajaron a las sombras del Hades,\\
que Aquiles, el de raudos pies, los dejó difuntos\\
entre los chuecos bueyes y las blancas ovejas.\\
Mi madre que a la falda de Placo nemoroso\\
reinaba también, trájola aquél, sordo a sus quejas,\\
con el botín. Mas, luego que mediante un cuantioso\\
rescate, concediónos que libertada fuera,\\
la inmoló en el alcázar Artemis la flechera.\\
Hoy, Héctor, para mí eres padre y madre honorable,\\
y hermano, a más de esposo, como ninguno amable.\\
Ten piedad, y no dejes la torre sin tu ayuda.\\
No hagas a tu hijo huérfano y a tu mujer viuda.\\
Detén a los soldados cerca del cabrahigo,\\
que ahí la ciudad es débil y el bastión inseguro.\\
Ya los héroes más bravos que cuenta el enemigo,\\
tres veces intentaron allá expugnar el muro,\\
bajo los dos Ayaces, el claro Idomeneo,\\
los Atridas y el fuerte vástago de Tideo.\\
O alguien quen los oráculos sabe se lo ha indicado\\
o bien su propio arrojo los ha precipitado.\\
Respondióle el Grande Héctor del casco tremolante:\par
--Mujer, bien sé todo eso, pero tambén me aterra\\
lo que dirán, si tímido me aparto de la guerra,\\
troyanos y troyanas de peplo rozagante.\\
Ni está en mi mano hacerlo\footnote{En la impresión «hacrelo», otro error.}, que siempre en la porfía\\
fuí valeroso y supe mantenerme adelante,\\
sosteniendo la gloria de mi padre y la mía.\\
Bien lo tengo en mi mente y en mi pecho arraigado:\\
día vendrá en que, al último, perezca Ilión la santa,\\
y Príamo y su pueblo, con el buen fresno armado.\\
Mas, ni el dolor de Troya desvalida me espanta,\\
ni el de la misma Hécuba, ni el de Príamo augusto,\\
ni el de mis muchos bravos hermanos, que el desastre\\
dará por tierra como tu dolor, cuando entonces,\\
uno de esos aqueos tunicados de bronce,\\
sollozante y privada de libertad te arrastre.\\
O cuando en Argos tejas para otro, y compelida\\
a ir por agua a las fuentes Meseida o Hiperea,\\
en ti sea la dura fatalidad cumplida.\\
Y quizá exclame alguno que llorando te vea:\\
esta es la esposa de Héctor que fué el primer guerrero\\
entre aquellos troyanos que ante Ilión combatieron.\\
Así dirán, y nuevo dolor ha de angustiarte\\
de haber perdido al hombre capaz de liberarte.\\
Pero antes cubra el polvo mis mortales despojos,\\
que oiga tu grito o vea tu rapto con mis ojos.\\
Dice y tiende sus brazos al niño; mas, gritando,\\
estréchase éste al seno de la apuesta nodriza,\\
pues lo asusta su tierno padre al irse allegando,\\
con el bronce y la mata de crin que tremolando\\
en el timbre del yermo, feramente se riza.\\
Padre y madre sonríen, y el bravo Héctor se quta\\
el centelleante casco que en tierra deposita.\\
Besa luego al infante y en sus brazos lo mece,\\
y a Zeus y demás dioses dirige así sus preces:\par
--Oh, Zeus y demás dioses, haced que mi hijo un día,\\
cual yo entre los troyanos, sea fuerte y glorioso.\\
Que sobre Ilión se vea reinando poderoso,\\
y en los comates digan de él por su bizarría:\\
``Pero éste vale mucho más que su padre!" Cuando\\
vuelva con los sangrientos despojos de un guerrero,\\
el alma de su buena madre regocijando.\par
Dice así y en los brazos de la querida esposa\\
depone al pequeñuelo, que ella en su perfumado\\
seno acoge, sonriendo todavía llorosa.\\
Entonces él, notándolo, la acaricia apiadado,\\
\dlgo{y dícele:}{--Oh, mi triste, no temas por mí, o creas}\\
que alguien me arroje\footnote{En el impreso «arorje».} al Hades, contrariando el destino,\\
pues no hay mortal que, tímido o bravo, eluda el sino\\
con que nació. Ya a casa vuelve y a tus tareas\\
de telar y del huso, ya  ver que las criadas\\
hagan su parte, mientras del trabajo guerrero\\
se encargan los soldados de Ilión, y yo el primero.\par
\pend
\endnumbering